\section{Setting of the problem}

In this chapter, we study an ML-inspired approach to goodness-of-fit testing known as
\textbf{Classifier-Accuracy Testing} (CAT), first explicitly described by \cite{friedmanCAT}
and later studied in a rigorous minimax framework by \cite{gerberclas}. The central
idea is to reduce a distributional testing problem to a classification problem: one trains a binary
classifier on synthetic labeled samples, then uses its predictions on the actual data as a test
statistic.

\begin{definition}[Total variation distance]
Let $\mu, \nu$ be two probability measures on $(\mathcal{X}, \mathcal{A})$. Their
\emph{total variation distance} is:
\[
    \mathrm{TV}(\mu, \nu) = \sup_{E \subset \mathcal{X}} |P(E) - Q(E)|
    = \frac{1}{2}\int_{\mathcal{X}} |f - g|\,d\lambda,
\]
where the second equality holds when $\mu, \nu$ admit densities $f, g$ w.r.t.\ a common
dominating measure $\lambda$.
\end{definition}

Given a sample $X$ of size $2n$ from an unknown $P_X$ and a sample $Y$ of size $2n$ from a
known null $P_Y$, we wish to test $H_0 : P_X = P_Y$ versus $H_1 : \mathrm{TV}(P_X, P_Y) \geq
\varepsilon$. A test $\psi : \mathcal{X} \to \{0,1\}$ achieves error at most $\delta \in
(0, \frac{1}{2})$ if $\max_{i}\max_{P \in H_i}\mathbb{P}(\psi(S) \neq i) \leq \delta$.

\medskip\noindent\textbf{The CAT framework.} Both samples are split as $X = X^{tr} \sqcup X^{te}$,
$Y = Y^{tr} \sqcup Y^{te}$. A classifier $C : \mathcal{X} \to \{0,1\}$ is trained on
$D^{tr} = \{(X_i, 1)\} \cup \{(Y_i, 0)\}$, and represented through its
\textbf{separating set} $S := C^{-1}(\{1\})$.

\begin{definition}[Classifier-accuracy statistic]
Given datasets $\{A_i\}_{i=1}^{a}$, $\{B_j\}_{j=1}^{b}$:
\[
    T_S(A, B) := \frac{1}{a}\sum_{i=1}^{a} \mathbf{1}_{A_i \in S}
               - \frac{1}{b}\sum_{j=1}^{b} \mathbf{1}_{B_j \in S}.
\]
\end{definition}

The test rejects $H_0$ when $|T_S(X^{te}, Y^{te})|$ exceeds a threshold. A key finding of
\cite{gerberclas} is that restricting $C$ to binary output does not sacrifice minimax optimality:
what matters is not classification accuracy, but finding an \emph{imbalanced} $S$ with large
separation $P_X(S) - P_Y(S)$ and small size $\tau(S) = \min\{P_X(S)P_X(S^c), P_Y(S)P_Y(S^c)\}$.

The main result of this chapter is the following lower bound on the sample complexity.

\begin{theorem}[{\cite[Table~1]{gerberclas}}]
\label{catcomplexitythm}
Consider the goodness-of-fit problem $H_0 : P_X = P_0$ versus $H_1 : \mathrm{TV}(P_X,P_0)
\geq \varepsilon$ with $P_0 \in \mathcal{P}_G(s, C)$. Then the minimax sample complexity satisfies:
\[
    n(\epsilon,\delta,\mathcal{P}_G(s,C_G))= n(\varepsilon, \delta, s)
    \;\gtrsim\;
    \frac{\sqrt{\log(1/\delta)}}{\varepsilon^{\frac{4s+1}{2s}}}
    + \frac{\log(1/\delta)}{\varepsilon^{2}}.
\]
\end{theorem}

The first term reflects the difficulty of nonparametric detection and depends on the smoothness $s$;
the second is a universal information-theoretic floor for binary testing at confidence $\delta$.
The proof is developed in the following sections.

\section{Preliminary results } 
The main difficulty of this kind of frameworks consists in finding a good separating set $S \subset \mx$, but we firstly need to define what "good" means. 
\begin{definition}
    Let $S \subseteq \mx$ be a measurable set, and $\mu,\nu$ be two probability measures on $\mx$. We define the separation and size of $S$ respectively as follows : 
    \begin{equation*}
        \sep(S) = \mu(S)-\nu(S) \ \text{ and } \ \tau(S) = \min(\mu(S)\mu(S^c),\nu(S)\nu(S^c))
    \end{equation*}
\end{definition}
\begin{lemma}\label{lem:lemma1}\citet[Lemma~1]{gerberclas}\\
    Consider the hypothesis testing problem $\H_0 : p = q$ versus any other alternative $\H_1$. Suppose that the learner has constructed a separating set $S$ such that : 
    \begin{equation*}
        \forall \mu,\nu \in \H_1 : \left|\sep(S)\right| = \left|\mu(S)-\nu(S)\right| \geq \underbar{$\sep$} \ \text{ and }\ \forall\mu,\nu \in \H_0 \cup \H_1 : \min(\mu(S)\mu(S^c),\nu(S)\nu(S^c)) \leq \bar{\tau}
    \end{equation*}
    Then, using only the knowledge of $\bar{\tau}$, the classifier-accuracy test with samples of size $n$ from both $p$ and $q$ and an appropriate threshold achieves type-I and type-II errors at most $\delta$ provided that : 
    \begin{equation*}
        n \geq c\frac{\log(\frac{1}{\delta})}{\underbar{$\sep$}}\Bigl(1+ \frac{\bar{\tau}}{\underbar{$\sep$}}\Bigr)
    \end{equation*}
    for a large enough constant $c >0$
\end{lemma}
\begin{remark}
With \Cref{lem:lemma1}, one can imagine that a separating set can be referred to as "good" if $\left|\sep(S)\right|$ is large under $H_1$, and $\tau(S)$ is small under both $H_0$ and $\H_1$ with probability $1-\delta$
\end{remark}
In order to prove \cref{lem:lemma1}, we will need a prior result on the Bernstein concentration of the classifier accuracy statistics. 
\begin{lemma}\label{lem:lemma2} \citet[Lemma~8]{gerberclas}\\
Suppose $A_1,A_2,\dots,A_n \overset{iid}{\sim}\text{Ber}(p)$ and $B_1,B_2,\dots,B_m \overset{iid}{\sim} \text{Ber}(q)$. Let $\tau = \min(p(1-p),q(1-q))$, and let $\bar{A},\bar{B}$ be their empirical means, respectively. There exists a universal constant $c>0$ such that :
\begin{equation*}
    \begin{aligned}
        \P\Bigg(\left|\bar{A}-\bar{B}\right|\leq \frac{1}{2}\left|p-q\right| -\frac{1}{2}\sqrt{\frac{c\log(1/\delta)\tau}{\min(n,m)}} -\frac{1}{2}\frac{c\log(1/\delta)}{\min(n,m)} \Bigg) &\leq \delta, \\
        \P\Bigg(\left|\bar{A}-\bar{B}\right|\geq 2\left|p-q\right| + 2\sqrt{\frac{c\log(1/\delta)\tau}{\min(n,m)}} + 2\frac{c\log(1/\delta)}{\min(n,m)} \Bigg) &\leq \delta
    \end{aligned}
\end{equation*} 
\end{lemma}
\begin{proof}[Proof of \cref{lem:lemma1}]
    Using $n$ test samples from $(X,Y)$, consider the following classifier-accuracy test: we accept $\H_0$ if  
    \begin{equation*}
        \Bigg| \frac{1}{n} \sum_{i=1}^{n} \ind{X_i \in S}-\ind{Y_i \in S} \Bigg| \leq \sqrt{\frac{c\bar{\tau}\log(1/\delta)}{n}} + \frac{c\log(1/\delta)}{n}
    \end{equation*}
    and reject the null hypothesis otherwise. Above, $c>0$ is a large absolute constant, and note that the threshold only relies on the knowledge of $\bar{\tau},n$ and $\delta$\\
    To analyse type-I and type-II errors , first, assume that $\H_0$ holds. Since $\sep(S) = 0$ under the latter hypothesis, the second statement of \Cref{lem:lemma2} implies that we accept $\H_0$ with probability at least $1-\frac{\delta}{2}$ if $c$ is large enough.
    Moreover, if $\H_1$ holds, with probability at least $1-\frac{\delta}{2}$, the first statement of \Cref{lem:lemma2} implies that : 
    \begin{equation*}
        \Bigg| \frac{1}{n}\sum_{i=1}^{n} \ind{X_i \in S} -\ind{Y_i \in S} \Bigg| \geq \left| \underbar{$\sep$} \right| + \Bigl(\sqrt{\frac{c\bar{\tau}\log(1/\delta)}{n}} + \frac{c\log(1/\delta)}{n}\Bigr)
     \end{equation*}
     By the assumed lower bound $n \geq c\frac{\log(\frac{1}{\delta})}{\underbar{$\sep$}}\Bigl(1+ \frac{\bar{\tau}}{\underbar{$\sep$}}\Bigr)$, we will reject $\H_0$ as desired.
\end{proof}
\begin{prop}\label{prop1}\citet[Prop~1]{gerberclas}
    Let $\P$ and $\Q$ be two probability measures on a probability space $\mx$, and let $\mc$ be a possibly randomized classifier. Then the following inequality is true : 
    \begin{equation*}
        TV(\P,\Q) \leq 2\sup \bigl\{\P(\mc(X) =0) -\Q(\mc(X)=0) \text{ over } X \ \text{satisfying} \ \P(\mc(X)=0) = \Q(\mc(X)=1) \bigr\}
    \end{equation*}
    In the inequality above, the constant 2 is tight.
\end{prop}
In order to prove \Cref{prop1} we will need the following lemma : 
\begin{lemma}\label{sixseven}
    Let $\mu$ be a non-negative measure on a measurable space $\mx$, and let $a,b:\mx\rightarrow\R_+$ such that $\int a(x) d\mu(x) >0$ and $\forall x \in \mx : b(x) = 0 \implies a(x) = 0$. Then the following holds : 
    \begin{equation*}
        \inf_{x \in \supp(\mu)}\frac{a(x)}{b(x)} \leq \frac{\int a(x) d\mu(x)}{\int b(x) d\mu(x)} \leq \sup_{x \in \supp(\mu)}\frac{a(x)}{b(x)}
    \end{equation*}
\end{lemma}
\begin{proof}[Proof of \Cref{prop1}]
    Let $p,q$ be the densities of $\P,\Q$ with respect to a common dominated measure (e.g. $\P+\Q$), and define $E \coloneqq \left\{p(x) > q(x)\right\}$, then $TV(\P,\Q) = \P(E)-\Q(E) \geq 0$. Without loss of generality, assume that $\P(E)+Q(\E) \geq 1$, and for $t \in [0,1]$, define : \begin{equation*}
        E_t \coloneqq \Bigl\{.\frac{p(x)-q(x)}{p(x)+q(x)} \geq t \Bigr\}
    \end{equation*}
    Then, the map $t \mapsto \P(E_t)+\Q(E_t)$ is non-increasing and left-continuous, and note that $E_0 = E$ and $E_1 = \emptyset$.
    Hence, the maximum $t^* \coloneqq \max \bigl\{t \in [0,1] : \P(E_t) +\Q(E_t) \geq 1\bigr\}$ exists. Now, define the randomized classifier $\mathcal{C}$ as follows : 
    \begin{equation*}
        \mathcal{C}(x) \coloneqq \begin{cases*}
            0 \ &if $x \in E_{(t^*)^+}$ \\
            1  \ &if $x \notin E_{t^*}$ \\
            \text{Ber}(x) &if $x \in E_{t^*}\backslash E_{(t^*)^+}$
        \end{cases*}
    \end{equation*}
    Where : \begin{equation*}
        E_{(t^*)^+} \coloneqq \cap_{t>t^*}E_t \ \text{ and } \ r = \frac{1-\P(E_{(t^*)^+})-\Q(E_{(t^*)^+})}{\P(E_{t^*})+\Q(E_{t^*})-\P(E_{(t^*)^+})-\Q(E_{(t^*)^+})} \in [0,1]
    \end{equation*}
    Note that this classifier is balanced since 
        $\P(\mathcal{C}(X)=0)+ \Q(\mathcal{C}(X) = 0) = 1$ .\\
    Now, for $t \in [0,1]$, define : 
    \begin{equation*}
        f(t) \coloneqq \begin{cases*}
            \frac{\P(E_t)-\Q(E_t)}{\P(E_t)+\Q(E_t)} & if $\P(E_t) + \Q(E_t) >0$ \\ 
            1 & otherwise 
        \end{cases*}
    \end{equation*}
    Now, let $0\leq t \leq s \leq 1$, we want to show that $f(t) \leq f(s)$. Without loss of generality, assume that $f(s) < 1$ and that $\P(E_t) + \Q(E_t) >0$. Notice also that (it can be proven using simple algebraic manipulations): 
    \begin{equation*}
        f(t) \leq f(s) \Leftrightarrow \frac{\int_{E_t \backslash E_s} p-q}{\int_{E_t \backslash E_s} p+q} \leq \frac{\int_{E_s} p-q}{\int_{E_s} p+q}
    \end{equation*}
    However, this inequality follows from \Cref{sixseven}. 
    Thus, it holds that : \begin{equation*}
        \frac{\P\bigl(\mathcal{C}(X) =0\bigr) - \Q\bigl(\mathcal{C}(X) = 0\bigr)}{\P\bigl(\mathcal{C}(X) =0\bigr) + \Q\bigl(\mathcal{C}(X) = 0\bigr)} \geq f(t^*)\geq f(0) = \frac{\P(E)-\Q(E)}{\P(E)+\Q(E)}
    \end{equation*}
    Plugging in $\P\bigl(\mathcal{C}(X) =0\bigr) + \Q\bigl(\mathcal{C}(X) = 0\bigr) = 1 $ and $\P(E)+\Q(E)\leq 2$ yields the final result. 
    For the tightness, consider $p(x) = \ind{[0,1],q(x) = (1+\epsilon)\ind{[0,\frac{1}{1+\epsilon}]}, \mathcal{C}(x) = \ind{(\frac{1}{2+\epsilon},1]}}$ and make $\epsilon \rightarrow 0$.
\end{proof}


\section{Performance of the CAT test in the Gaussian case}
Suppose one has two samples $X,Y$ of size $n\in \N$ from $\mu_{\theta^X} \coloneqq \bigotimes_{j=1}^\infty \mathcal{N}(\theta_j^X,1)$ and $\mut{\theta^Y}$, respectively and such that $\mu^X,\mu^Y$ have bounded Sobolev norm. 
Moreover, assume that $TV(\mut{\theta^X},\mut{\theta^Y}) \geq \epsilon >0$, and let $\hat{\theta}^X, \hat{\theta}^Y$ denote the empirical means from the two samples respectively.
Define the following statistic : 
\begin{equation}\label{statseparating}
    T(Z) = 2\sum_{j=1}^{J} \Bigl(\hat{\theta}_j^X-\hat{\theta}_j^Y\Bigr)\Biggl(Z_j - \frac{\hat{\theta}_j^X+\hat{\theta}_j^Y}{2}\Biggr)
\end{equation}
for some $J\in \N$ to be specified. We construct the separating set as follows : \begin{equation}\label{sepset}
    \hat{S} \coloneqq \bigl\{Z \in \R^\N : T(Z)\geq 0\bigr\}
\end{equation}
The performance of this construction of separating set is summarized in the statement of the following result : 
\begin{prop}\citet{gerberclas}\label{CATperf}
    There exists universal constants $c,c'$ such that for $J = \big\lfloor c\epsilon^{-\frac{1}{s}}\big\rfloor$ the following inequality holds provided $n \gtrsim \frac{1}{c'}n(\epsilon,\delta,s)$ : 
    \begin{equation}\label{ineqperf}
        \P\Bigg( \mut{\theta^X}(\hat{S}) - \mut{\theta^Y}(\hat{S}) \geq c'\cdot \min\Big(\sqrt{n\epsilon^{\frac{1}{s}}},\frac{1}{\epsilon}\Big)\cdot \epsilon^2 \Bigg) \geq 1 -\delta
    \end{equation}
\end{prop}
In order to prove this proposition, we will need few lemmas and prior results. 
\begin{lemma}[Lipschitz concentration of Gaussians]\label{LipschitzConc}
    Let $Q$ be a $d$-dimensional standard gaussian, and let $f : \R^d \rightarrow \R$ be $L$-Lipschitz. Then $f(Q)$ is sub-Gaussian with variance proxy $L^2$
\end{lemma}
\begin{proof}
    See \hyperref[appendix]{Appendix}
\end{proof}
\begin{lemma}\label{gaussianexpec}
    Let $a,b\in \R$, and let $Z$ be a standard normal. Then : \begin{equation*}
        \E \Phi(aZ+b) = \Phi\bigg(\frac{b}{\sqrt{1+a^2}}\bigg)
    \end{equation*}
\end{lemma}
\begin{proof}
    See \hyperref[appendix]{Appendix}
\end{proof}
\begin{lemma}\label{lemma13gerber}
    Let $\sep(\hat{S}) = \mut{\theta^X}(\hat{S})- \mut{\theta^Y}(\hat{S})$ as defined above. There exists universal constants $c_i >0$, i = 1,2,$\cdots$,5 such that for $J = \big\lfloor c_1\epsilon^{-\frac{1}{s}}\big\rfloor$ and for $t \geq 0$, one has : 
    \begin{align}
        \E[\sep(\hat{S})] + \frac{c_2}{\sqrt{n}} &\geq \frac{c_3 \epsilon^2}{\epsilon + \sqrt{\frac{J}{n}}} \\
        \P\Biggl(\Big|\sep(\hat{S})-\E \sep(\hat{S}) \Big| \geq t+ \frac{c_4}{\sqrt{n}}\Biggr) &\leq 2 \exp(-c_5nt^2)
    \end{align}
\end{lemma}
\begin{proof}[Proof of \Cref{lemma13gerber}]
    Let $\norm{\cdot}, \scal{\cdot}{\cdot}$ for the $\l^2$ norm/inner product restricted to the first $J$ coordinates. Notice that given $\hat{\theta}^X,\hat{\theta}^Y$, $T(\theta)$ is a Gaussian random variable with : 
    \begin{equation*}
        \E T(\theta) = \normp[2]{\hat{\theta}^Y - \theta} - \normp[2]{\hat{\theta}^X-\theta}  \ \text{ and } \ \var{T} = 4 \normp[2]{\hat{\theta}^X-\hat{\theta}^Y}
    \end{equation*}
    Now, define the random vectors $U,V$ as follows : 
    \begin{align*}
        U &= \bigl\{ \hat{\theta}^X_j - \hat{\theta}^Y_j\bigr\}_{j=1}^J \\
        V &= \bigl\{\hat{\theta}^X_j + \hat{\theta}^Y_j\bigr\}_{j=1}^J
    \end{align*}
    One can easily see that these two vectors are independent, Gaussians with covariance matrix $\frac{2}{n}I_J$ and means equal to the first $J$ coordinates of $\theta^X \minusplus \theta^Y$ respectively. Let $\Phi$ be the cumulative distribution function of the standard Gaussian and $\phi$ be its density.
    It can be shown that the separation can be written as : \begin{equation*}
        \sep(\hat{S}) = f(\hat{\theta}^X)- f(\hat{\theta}^Y)
    \end{equation*}
    Where $f$ is defined as follows : 
    \begin{equation*}
        f(\theta) = \Phi\Bigg( \frac{\normp[2]{\hat{\theta}^Y -\theta}-\normp[2]{\hat{\theta}^X -\theta}}{2\normp[2]{\hat{\theta}^X -\hat{\theta}^Y}} \Bigg) = \Phi\Bigg(-\frac{1}{2}\bigscal{V}{\frac{U}{\norm{U}}} + \bigscal{\theta}{\frac{U}{\norm{U}}}\Bigg)
    \end{equation*}
    In order to make the dependences more clear, write $g(U,V) = f(\theta^X)-f(\theta^Y)$. Now, differentiating $g$ and using the fact that $\phi$ is $\frac{1}{\sqrt{2\pi e}}$-Lipschitz yields : 
    \begin{align*}
        \norm{\nabla_V g(U,V)}  &= \begin{aligned}[t]
            \bignorm{&-\frac{1}{2}\frac{U}{\norm{U}}\Bigg(\phi\Bigl( -\frac{1}{2}\bigscal{V}{\frac{U}{\norm{U}}} + \bigscal{\theta^X}{\frac{U}{\norm{U}}} \Bigr)\\
            &- \phi\Bigl(-\frac{1}{2}\bigscal{V}{\frac{U}{\norm{U}}} + \bigscal{\theta^Y}{\frac{U}{\norm{U}}} \Bigr) \Bigg)}
        \end{aligned}\\
        & \leq \frac{1}{\sqrt{8\pi e}} \Big|\bigscal{\theta^X-\theta^Y}{\frac{U}{\norm{U}}} \Big| \\ 
        & \leq \frac{C_G}{\sqrt{8 \pi e}}
    \end{align*}
    By Lipschitz concentration of Gaussian random variable (\Cref{LipschitzConc}), one can conclude that $g - \E[g|U]$ is sub-Gaussian with variance proxy $\frac{C_G^2}{4\pi en}$. Now, we study the concentration of $\E[g|U]$. First, note that : \begin{equation*}
        -\frac{1}{2} \bigscal{V}{\frac{U}{\norm{U}}} + \bigscal{\theta}{\frac{U}{\norm{U}}} \Bigg| U \sim \mathcal{N}\Biggl(\bigscal{\theta-\frac{1}{2}(\theta^X+\theta^Y)}{\frac{U}{\norm{U}}},\frac{1}{2n} \Biggr)
    \end{equation*}
    Then, using the independence of $U,V$ and \cref{gaussianexpec}, we obtain : \begin{align*}
        \E[g(U,V)|U] &= \E[f(\theta^X)-f(\theta^Y)|U] \\ 
        &= \Phi\Biggl(\frac{W}{\sqrt{4+\frac{2}{n}}}\Biggr)-\Phi\Biggl(-\frac{W}{\sqrt{4+\frac{2}{n}}}\Biggr)
    \end{align*}
    Where $W \coloneqq \scal{\theta^X-\theta^Y}{\frac{U}{\norm{U}}}$. To simplify notations, let $\varphi \coloneqq \Phi\Bigl(\frac{\cdot}{\sqrt{4+\frac{2}{n}}}\Bigr)$, by Lipschitness of $\Phi$, we obtain that for $t \geq 0$: \begin{align*}
        \P\Bigl( \Big|\varphi(W)-\E\varphi(W) \Big| \geq t\Bigr) &\leq \P\Bigl(\Big|\varphi(W) - \varphi(\E W)\Big|\geq t-\lipnorm{\varphi}\sqrt{\var{W}}\Bigr) \\
        &\leq  \P \Bigl(|W-\E W| \geq \frac{t}{\lipnorm{\varphi}}-\sqrt{\var{W}} \Bigr)
    \end{align*}
    Moreover, one can get an analogous inequality for $-W$. Now, the last ingredient is to show that $W$ concentrates well. 
\begin{lemma}\label{youssef}
    W is a sub-Gaussian random variable with variance proxy $\frac{8}{n}$
\end{lemma}

We have already seen that : 
\begin{equation*}
    \E[\sep(\hat{S})] = \E[\varphi(W)-\varphi(-W)]
\end{equation*}
Again, by Lipschitzness and \Cref{gaussianexpec}, we have that $|\E [\varphi(W)]-\varphi(\E W)|\lesssim \frac{1}{\sqrt{n}}$, hence : 
\begin{equation}\label{omega}
    \varphi(\E W)- \varphi(-\E W) \leq \E[\sep(\hat{S})] + \Omega(\frac{1}{\sqrt{n}})
\end{equation}
Where the implied constant is universal. In order to simplify notations, let $\tau \coloneqq \theta^X-\theta^Y, \sigma^2 = \frac{2}{n}$ and let $Q$ be a $\mathcal{N}(0,I_J)$ random vector, then : 
\begin{equation*}
    \E W = \E \bigg[ \bigscal{\tau}{\frac{\tau+\sigma Q}{\norm{\tau+\sigma Q}}}\bigg] = \frac{1}{\sigma}\E[\scal{\tau}{\nabla_Q \norm{\tau+\sigma Q}}] = \frac{1}{\sigma}\E[\scal{\tau}{Q}\norm{\tau+\sigma Q}]
\end{equation*}
Where the last equality is given by Stein's identity. By rotational invariance of Gaussian distribution, we get : 
\begin{align*}
    \E W &= \frac{\norm{\tau}}{\sigma}\E\Bigg[Q_1 \sqrt{(\norm{\tau}+\sigma Q_1)^2+\sigma^2Q_2^2 + \cdots \sigma^2Q_J^2}\Bigg]\\ 
    &= \frac{\norm{\tau}}{\sigma}\E\Bigg[Q_1 \sqrt{(\norm{\tau}+\sigma Q_1)^2+\sigma^2Q_2^2 + \cdots \sigma^2Q_J^2} - Q_1 \sqrt{\norm{\tau}^2+\sigma^2 Q_1^2+\sigma^2 Q_2^2 + \cdots \sigma^2Q_J^2}\Bigg]\\
    &= 2\norm{\tau}^2 \E \Bigg[ \frac{Q_1^2}{Q_1 \sqrt{(\norm{\tau}+\sigma Q_1)^2+\sigma^2Q_2^2 + \cdots \sigma^2Q_J^2} + Q_1 \sqrt{\norm{\tau}^2+\sigma^2 Q_1^2+\sigma^2 Q_2^2 + \cdots \sigma^2Q_J^2}} \Bigg]
\end{align*}
Now, by Cauchy-Schwartz inequality, and the fact that $\E |Q_1| = \frac{2}{\pi}$ is an absolute constant, we get that : 
\begin{equation*}
    (\E |Q_1^2|)^2 \lesssim \E \Bigg[ \frac{Q_1^2}{\sqrt{(\norm{\tau}+\sigma Q_1)^2+\sigma^2Q_2^2 + \cdots \sigma^2Q_J^2} + \sqrt{\norm{\tau}^2+\sigma^2 Q_1^2+\sigma^2 Q_2^2 + \cdots \sigma^2Q_J^2}} \Bigg] \times (\norm{\tau}+\sigma \sqrt{J})
\end{equation*}
Plugging everything into our expression for $\E W$, this yields : 
\begin{equation*}
    \E W \gtrsim \frac{\normp[2]{\tau}}{\norm{\tau} + \sigma\sqrt{J}}
\end{equation*}
To clarify notation, let us now write $\norm{\cdot}_J$ for the $\l_2$-norm restricted to
the first $J$ coordinates. Taking $J = c \epsilon^{-1/s}$ it holds that
\[
  \norm{\tau}_J^2 \;=\; \norm{\tau}^2 - \sum_{j > J}\tau_j^2
   \;\geq\; \norm{\tau}^2 - J^{-2s}\sum_{j > J}\tau_j^2 j^{2s}
   \;=\; \norm{\tau}^2 - c^{-2s}\epsilon^2 \norm{\tau}_s^2.
\]
Since $\norm{\tau}_s \lesssim 1$ and $\norm{\tau} \geq \epsilon$ by assumption, we see
that for large enough universal constant $c$ we have $\norm{\tau}_J \geq \epsilon/2$.
Since the map $x \mapsto x^2/(x + a)$ is increasing for $x, a > 0$, it follows that
\[
  \E W \;\gtrsim\; \frac{\epsilon^2}{\epsilon + \sqrt{J/n}}
\]
for a universal implied constant (using $\sigma\sqrt{J} = \sqrt{2J/n}\asymp\sqrt{J/n}$).
By the inequality $\Phi(x) - \Phi(-x) \geq x/2$ for $x \in [0,1]$ we obtain
\[
  \varphi(\E W) - \varphi(-\E W) \;\geq\; 1 \wedge \E W/2,
\]
which, combined with \Cref{omega}, completes the proof.

\end{proof}
\begin{proof}[Proof of \Cref{CATperf}]
    Apply \Cref{lemma13gerber} and set : 
    \begin{equation*}
        c = c_1 \ \text{ and } \ M \coloneqq \frac{c_3 \epsilon^2}{\epsilon+ \sqrt{\frac{J}{n}}}
    \end{equation*}
    By \Cref{lemma13gerber}, $\E[\sep(\hat{S})] \geq M -\frac{c_2}{\sqrt{n}}$. Now, choose $t = \sqrt{\frac{\log(2/\delta)}{c_5 n}}$ so that $2\exp(-c_5nt^2) = \delta$. By the tail bound given by \Cref{lemma13gerber} : 
    \begin{equation*}
        \P \Bigl(\sep(\hat{S}) <  \E[\sep(\hat{S})]-t-\frac{c_4}{\sqrt{n}}\Bigr) \leq \delta
    \end{equation*}
    Hence, with probability at least $1-\delta$, 
    \begin{equation}\label{ineqsep}
        \sep(\hat{S}) \geq M -\frac{c_2+c_4}{\sqrt{n}}-\sqrt{\frac{\log(2/\delta)}{c_5 n}} \geq M -C\sqrt{\frac{\log(2/\delta)}{n}}
    \end{equation}
    for some universal constant $C>0$. Now, for $J = \lfloor c\epsilon^{-\frac{1}{s}}\rfloor$, we have $\frac{J}{n} \asymp \frac{\epsilon^{-\frac{1}{s}}}{n}$, hence : 
    \begin{equation*}
        \epsilon+ \sqrt{\frac{J}{n}} \asymp \begin{cases*}
            \epsilon &\text{if} $\sqrt{n} \epsilon^{1+\frac{1}{2s}} \geq 1$ \\ 
            \sqrt{\frac{J}{n}} & ,otherwise
        \end{cases*}
    \end{equation*}
    Using the asymptotics above, one can easily check that : 
    \begin{equation*}
        M \asymp \Bigl(\sqrt{n}\epsilon^{\frac{1}{s}}\wedge \frac{1}{\epsilon}\Bigr)\epsilon^2
    \end{equation*}
    Also, the second term in \Cref{ineqsep} is of order $\sqrt{\frac{\log(1/\delta)}{n}}$, under the condition $n \gtrsim \frac{1}{c'}n(\epsilon,\delta,s)$ (where $c' = \frac{c_3}{2}$ after absorbing all universal constants), 
    this term is at most $\frac{M}{2}$. Therefore, with probability at least $1-\delta$ : 
    \begin{equation*}
        \sep(\hat{S}) \geq \frac{M}{2} \asymp \Bigl(\sqrt{n}\epsilon^{\frac{1}{s}}\wedge \frac{1}{\epsilon}\Bigr)\epsilon^2
    \end{equation*}
    Which is the statement of \Cref{CATperf}. (Note that one can obtain the last inequality by adjusting the constants)
\end{proof}
\section{Sample complexity of CAT for Gaussian sequence model}
Given two hypotheses $H_0,H_1$, our goal is to find a lower bound for the minimum achievable worst-case error. 
Firstly, we use the fact that for any random probability distributions $P_0,P_1$ with $\P(P_i\in H_i)=1$ for $i=1,2$, one has that : 
\begin{equation}\label{lowerbounderror}
    \min_{\psi}\max_{i=0,1}\sup_{P\in H_1} \P_{S\sim P} \big(\psi(S) \neq i \big) \geq \frac{1}{2}\big(1-TV(\E P_0,\E P_1)\big)
\end{equation}
Where $\E P$ denotes the mixture of $P$. Also, note that deriving a lower bound of order $\delta$ on the minimax error is equivalent to finding mixtures $\E P_0,\E P_1$ such that $1-TV(\E P_0,\E P_1)$ is of order $\delta$ (is $\Omega(\delta))$

\begin{lemma}\label{rubignolle}\cite{tsybakov}
    Let $\mu,\nu$ be two probability measures, then : 
    \begin{equation*}
        1-TV(\mu,\nu) \geq \frac{1}{2} e^{-KL(\mu \Vert \nu)} 
    \end{equation*}
    Where $KL(\cdot \Vert \cdot)$ denote the Kullback-Leibler-divergences respectively. 
\end{lemma}
\begin{proof} [Proof of \Cref{rubignolle}] We will show that each one of the two inequalities hold true using Jensen inequality. 
    Instead of the inequality above, one will show that : \begin{equation*}
        TV(\mu,\nu) \leq \sqrt{1-e^{-KL(\mu\Vert\nu)}} 
    \end{equation*}
    First, notice that : \begin{equation*}
        TV(\mu, \nu) = 1- \int \bigl(d\mu \land d\nu\bigr) = \int \bigl(d\mu \lor d\nu\bigr) -1 
    \end{equation*}
    This implies that : \begin{align*}
        1 - TV(\mu,\nu)^2 &= \bigl(1 - TV(\mu,\nu)\bigr)\bigl(1+TV(\mu,\nu)\bigr) \\
        &= \int \bigl(d\mu \land d\nu\bigr)\int \bigl(d\mu \lor d\nu\bigr) \\
        &\stackrel{\text{C-S}}{\geq} \Biggl(\int \sqrt{fg}\Biggr)^2 \  \text{ (Where $f,g$ are the densities)}\\ 
        &= \exp\Biggl(2\log\Bigl(\int \sqrt{\frac{d\nu}{d\mu}} d\mu\Bigr)\Biggr) \\ 
        &= \exp\Biggl(2\log\E_\mu \Biggl(\sqrt{\frac{d\mu}{d\nu}}\Biggr)^{-1} \Biggr) \\
        &\geq \exp \Biggl(-\E_\mu \log\biggl(\frac{d\mu}{d\nu}\biggr)\Biggr) = \exp\biggl(-KL(\mu\Vert\nu)\biggr)
    \end{align*}
    And finally, one can conclude using the following inequality : $\sqrt{1-x} \leq 1-\frac{x}{2}$.
\end{proof}

Now that we have all the necessary tools, we will prove the main theorem of this chapter. We will split this proof into many parts : 
\begin{prop}\label{param}\cite{gerberclas}
    There exists a constant $K_0$ such that for every $\epsilon$ small enough, \begin{equation}\label{parameq}
        n(\epsilon,\delta,s) \geq K_0 \frac{log(1/\delta)}{\epsilon^2}
    \end{equation}
\end{prop}
\begin{proof}[Proof of \Cref{param}] 
    For a vector $\eta \in \{-1,1\}^\N$, define the following probability measure : 
    \begin{equation}\label{randomproba}
        \P_\eta \coloneqq \bigotimes_{j=1}^J \mathcal{N}(\eta_jc_1 \epsilon^{\frac{2s+1}{2s}},1) \otimes \bigotimes_{j>J} \mathcal{N}(0,1)
    \end{equation}
    Where $J = \big\lfloor c_2\varepsilon^{-\frac{1}{s}}\big\rfloor$. And let $\eta_1,\eta_2,\cdots$ be iid uniform in $\{-1,1\}$ and $\gamma_\eta$ be the mean vector of $\P_\eta$. One can see that for any such $\eta$ : 
    \begin{align*}
        \norm{\gamma_\eta}_s^2 = \sum_{j=1}^{\infty} j^{2s}\gamma_{\eta,j}^2 = \sum_{j=1}^{J}j^{2s}c^2\epsilon^{\frac{2s+1}{s}} &\leq c_1^2 \epsilon^{\frac{2s+1}{s}} \bigl( 2c_2\epsilon^{-\frac{1}{s}} \bigr)^{2s+1} \asymp c_1^2c_2^{2s+1}\\
        \normp[2]{\gamma_\eta} &= \sum_{j=1}^{\infty} \gamma_{\eta,j}^2 \asymp c_1^2c_2\epsilon^2
    \end{align*}
    Then for $C_G >0$ (Size of Sobolev's ellipsoid) fixed, one can choose $c_1,c_2$ independently of $\epsilon$ such that $\P_0,\P_\eta \in \mathcal{P}_G(s,C_G)$ almost surely and $\norm{\gamma_\eta} = 10\epsilon$. 
    Then, for $\epsilon \leq \frac{1}{10}$, one has : 
    \begin{equation*}
        TV(\P_0,\P_\eta) = 2 \Phi\Biggl(\frac{\norm{\gamma_\eta}}{2}\Biggr) - 1 \geq \epsilon
    \end{equation*}
    Now, take $P_0 = \P_0^{\otimes2n}, P_1 = \P_1^{\otimes n} \otimes \P_0^{\otimes n}$, then : \begin{equation*}
        KL(P_0\Vert P_1) = n KL(\P_0 \Vert \P_1) = nc_2 \epsilon^{-\frac{1}{s}}\frac{(c_1 \epsilon^{\frac{2s+1}{2s}})^2}{2} = K n\epsilon^2
    \end{equation*}
    Using \Cref{rubignolle} and \Cref{lowerbounderror}, we get the minimax error is at least : \begin{equation*}
        \frac{1}{2}\bigl(1-TV(P_0,P_1) \bigr)\geq \frac{1}{4}\exp(-KL(P_0\Vert P_1)) = \frac{1}{4}\exp(-Kn\epsilon^2) 
    \end{equation*}
    For this error to be less than $\delta$, it is necessary to have : \begin{equation*}
        \frac{1}{4}\frac{1}{4}\exp(-Kn\epsilon^2) \leq \delta \Leftrightarrow n \geq \frac{\log(1/(4\delta))}{K\epsilon^2}
    \end{equation*}
    Hence $n \gtrsim \frac{\log(1/\delta)}{\epsilon^2}$ as claimed.
\end{proof}
\begin{prop}\label{nonparam}\cite{gerberclas}
    There exists a constant $K_1 >0$ such that for every $\epsilon$ small enough : 
    \begin{equation}\label{nonparameq}
        n(\epsilon,\delta,\mathcal{P}_G(s,C_G)) \geq K_1 \frac{\sqrt{\log(1/\delta)}}{\epsilon^{\frac{4s+1}{2s}}}
    \end{equation}
\end{prop}
\begin{proof}
    Let $\eta$ be a sequence of iid uniform signs (i.e. with components from $\{-1,1\}$ uniformly), and define the following probability measures : 
    \begin{equation*}
    P_0 = \P_0^{\otimes 2n},P_1 = \P_0^{\otimes n}\otimes \P_\eta^{\otimes n}
    \end{equation*} 
    By the independence of the components of $\eta$ and their identical distribution, for $\omega=c_1 \epsilon^{\frac{2s+1}{2s}}$ the following mixture can be written as follows : 
    \begin{equation*}
        \E\P_\eta^{\otimes n} = \bigotimes_{j=1}^J \Biggl(\frac{1}{2}\mathcal{N}(\omega,1)^{\otimes n}+\frac{1}{2}\mathcal{N}(-\omega,1)^{\otimes n}\Biggr) \otimes \bigotimes_{j>J}\mathcal{N}(0,1)^{\otimes n}
    \end{equation*}
    In order to simplify notations, define $M_n \coloneqq\frac{1}{2}\mathcal{N}(\omega,1)^{\otimes n}+\frac{1}{2}\mathcal{N}(-\omega,1)^{\otimes n}$, hence : 
    \begin{equation}\label{mixture}
        KL(\P_0^{\otimes n} \Vert \E\P_\eta^{\otimes n}) = J\cdot KL(\mathcal{N}(0,1)^{\otimes n}\Vert M_n)
    \end{equation}
\begin{lemma}\label{klmixture}
    For $\omega \in \R$ and $n \geq 1$ : 
    \begin{equation}
        KL\Biggl(\mathcal{N}(0,1)^{\otimes n}\bigg\Vert \frac{1}{2}\mathcal{N}(\omega,1)^{\otimes n} +\frac{1}{2}\mathcal{N}(-\omega,1)^{\otimes n} \Biggr) \leq \frac{n^2\omega^4}{4}
    \end{equation}
\end{lemma}
\begin{proof}[Proof of \Cref{klmixture}]
    See \hyperref[appendix]{Appendix}
\end{proof}
By combining \Cref{mixture} and \Cref{klmixture} and using the independence of the components, we obtain : 
\begin{equation*}
    KL(\P_0^{\otimes n}\Vert \E\P_\eta^{\otimes n}) \leq J\cdot \frac{n^2\omega^4}{4} = \frac{n^2}{4}\cdot J\omega^2 \cdot \omega^2
\end{equation*}
But $J\omega^2= \normp[2]{\gamma_\eta} = 100\epsilon^2$ and $\omega^2 = c_1^2 \epsilon^{\frac{2s+1}{s}}$ hence : 
\begin{equation*}
    KL(\P_0^{\otimes n}\Vert \E\P_\eta^{\otimes n}) \leq \frac{100c_1^2}{4}n^2 \epsilon^{\frac{4s+1}{s}} = \kappa \cdot n^2 \epsilon^{\frac{4s+1}{s}}
\end{equation*}
Now, by combining \Cref{lowerbounderror} and \Cref{rubignolle}, in order to get an error less than $\delta$, we need to have : 
\begin{equation*}
    \kappa n^2\epsilon^{\frac{4s+1}{s}} \geq \log\Bigl(\frac{1}{4\delta}\Bigr)
\end{equation*}
Which yields $n \gtrsim \frac{\sqrt{\log(1/\delta)}}{\epsilon^{\frac{4s+1}{2s}}}$
\end{proof}

\begin{proof}[Proof of \Cref{catcomplexitythm}]
    Using \Cref{param} and \Cref{nonparam}, we get that there exists two constants $K_0,K_1 >0$ such that \Cref{parameq} and \Cref{nonparameq} hold true. 
    Let $K = \min \{K_0,K_1\}$, one has : 
    \begin{align*}
        n &\geq \max\Biggl\{K_0 \frac{\log(1/\delta)}{\epsilon^2},K_1 \frac{\sqrt{\log(1/\delta)}}{\epsilon^{\frac{4s+1}{2s}}}\Biggr\}\\ 
        &\geq \max\Biggl\{K \frac{\log(1/\delta)}{\epsilon^2},K \frac{\sqrt{\log(1/\delta)}}{\epsilon^{\frac{4s+1}{2s}}}\Biggr\}\\ 
        &\geq \frac{K}{2}\Biggl(\frac{\log(1/\delta)}{\epsilon^2}+\frac{\sqrt{\log(1/\delta)}}{\epsilon^{\frac{4s+1}{2s}}}\Biggr)
    \end{align*}

\end{proof}